\documentclass{article}
\usepackage{svg}
\usepackage{amsmath}
\usepackage{geometry} % Required for customizing margins
\geometry{margin=1in} % Sets 1-inch margins on all sides

\title{ELCT 331 Control Systems Assignment 2}
\author{Clarence Hunter Frady}
\date{Fall 2023}
\begin{document}
\maketitle

\section*{Problem \#1[20 Points]}
Reduce the block diagram shown in the figure to unity feedback form and find the system characteristic equation. 
\newline
\begin{figure}[htbp]
  \centering
  \includesvg[width=\textwidth]{331_HW2_P1.svg}
\end{figure}

\newpage
\section*{Problem \#2[20 Points]}
The Figure shows the block diagram of a dc-motor control system. The signal \( N(s) \) denotes torque at the motor shaft.

\begin{enumerate}
    \item[a. ] Find the transfer function \( H(s) \) so that the output \( Y(s) \) is not affected by the disturbance torque \( N(s) \).
    \item[b. ] With \( H(s) \) as determined in part (a), find the value of \( K \) so that the steady-state value of \( e(t) \) is equal to 0.1 when the input is a unit-ramp function \( r(t)=tu_s(t) \), \( R(s)=\frac{1}{s^2} \), and \( N(s)=0 \). Apply the final-value theorem.
\end{enumerate}
\newline
\begin{figure}[htbp]
  \centering
  \includesvg[width=\textwidth]{331_HW2_P2.svg}
  \Large \( G(s) = \frac{K(s + 3)}{s(s + 1)(s + 2)} \)
\end{figure}


\newpage
\section*{Problem \#3[20 Points]}
\begin{enumerate}
    \item[a. ] Draw an equivalent Signal Flow Graph (SFG) for the system.
    \item[b. ] Find the following transfer functions by applying the gain formula of the SFG directly to the block diagram.
    \[
    \frac{Y(s)}{R(s)} \bigg|_{N=0} \quad \frac{Y(s)}{N(s)} \bigg|_{R=0} \quad \frac{E(s)}{R(s)} \bigg|_{N=0} \quad \frac{E(s)}{N(s)} \bigg|_{R=0}
    \]
    \item[c. ] Compare the answers by applying the gain formula to the equivalent SFG.
\end{enumerate}
\newline
\begin{figure}[htbp]
  \centering
  \includesvg[width=\textwidth]{331_HW2_P3.svg}
\end{figure}


\newpage
\section*{Problem \#4 [30 points]}

The figure shows the block diagram of a control system with conditional feedback. The transfer function \( G_p(s) \) denotes the controlled process, and \( G_c(s) \) and \( H(s) \) are the controller transfer functions.

\begin{enumerate}
    \item[a. ] Derive the transfer functions 
    \[
    \frac{Y(s)}{R(s)} \bigg|_{N=0}  \quad \frac{Y(s)}{N(s)} \bigg|_{R=0}
    \]
    Find \( \frac{Y(s)}{R(s)} \bigg|_{N=0} \) when \( G_c(s) = G_p(s) \).
    
    \item[b. ] Let
    \[
    G_p(s) = G_c(s) = \frac{100}{(s+1)(s+5)}
    \]
    Find the output response \( y(t) \) when \( N(s) = 0 \) and \( r(t) = u_s(t) \).

    \item[c. ] With \( G_p(s) \) and \( G_c(s) \) as given in part (b), select \( H(s) \) among the following choices such that when \( n(t) = u_s(t) \) and \( r(t) = 0 \), the steady-state of \( y(t) \) is equal to zero. (There may be more than one answer.)
\begin{center}
  \begin{minipage}{0.4\textwidth}
    \begin{align*}
      H(s) &= \frac{10}{s(s+1)} \\
      H(s) &= \frac{10(s+1)}{s+2}
    \end{align*}
  \end{minipage}
  \begin{minipage}{0.4\textwidth}
    \begin{align*}
      H(s) &= \frac{10}{(s+1)(s+2)} \\
      H(s) &= \frac{K}{s^n} \text{ (n is a positive integer. Select n.)}
    \end{align*}
  \end{minipage}
\end{center}

    
    
    Keep in mind that the poles of the closed-loop transfer function must all be in the left-half \( s \)-plane for the final-value theorem to be valid.
\end{enumerate}
\newline
\begin{figure}[htbp]
  \centering
  \includesvg[width=\textwidth]{331_HW2_P4.svg}
\end{figure}

\end{document}
